\bookchapter{Capítulo 4: Instalar e Gerenciar Apps — O “Pacote de Habilidades” do Smartphone}

\chapterintro{Os apps tornam o smartphone capaz de tudo, mas é preciso gerenciá-los bem. Assim como você organiza a estação de trabalho e atribui permissões a novos funcionários, instalar apps também exige um processo padronizado.}

\booksection{I. Lojas Oficiais: O “Supermercado de Apps” mais seguro}

\begin{itemize}
  \item \textbf{Android}: Google Play (no exterior) ou lojas próprias das marcas (como Huawei AppGallery, Xiaomi App Store).
  \item \textbf{iOS}: Apenas App Store.
  \item \textbf{HarmonyOS}: Huawei AppGallery.
\end{itemize}

Significado central de segurança: passam por revisão, atualizações garantidas, pagamento seguro. Guia de ação:\textbf{Priorize baixar apenas das lojas oficiais!}

\booksection{II. Modo Puro / Modo de Proteção: Ative o “Firewall” do telefone}

As principais marcas chinesas (Huawei, Xiaomi, OPPO, vivo, Honor) possuem integrado o\textbf{Modo Puro}ou\textbf{Modo de Proteção}, um interruptor de segurança no nível do sistema.\textbf{Recomenda-se fortemente que todos os usuários o ativem}.

\booksubsection{O que é exatamente?}

O Modo Puro (também chamado “Modo de Proteção”) tem como\textbf{função principal bloquear a instalação de apps de terceiros}: quando ativado, o sistema\textbf{permite apenas a instalação de apps da loja oficial}. Pacotes de instalação baixados de navegadores, QR codes, aplicativos de mensagens etc. são bloqueados diretamente, eliminando malware na origem.

\booksubsection{Público-alvo}

\begin{itemize}
  \item \textbf{Usuários idosos}: Não conseguem distinguir anúncios "oficiais" de "falsos". O Modo Puro é o amuleto mais seguro.
  \item \textbf{Usuários pouco sensíveis a permissões de apps}: Não precisam instalar apps de nicho, versões beta ou jogos de servidores estrangeiros.
  \item \textbf{Todos os usuários que desejam aumentar a segurança do dispositivo}: Mesmo para entusiastas, é melhor ativado do que desativado — se houver necessidade temporária, basta desativar momentaneamente.
\end{itemize}

\booksubsection{Caminho de ativação (exemplo Huawei/Xiaomi)}

\begin{itemize}
  \item \textbf{Huawei/HarmonyOS}: [Ajustes] > [Sistema e atualizações] >\textbf{[Modo Puro]}, selecione "Ativar".
  \item \textbf{Xiaomi}: [Ajustes] > [Privacidade] >\textbf{[Proteger privacidade]}ou [Ajustes] > [Configurações de apps] > [Instalação de apps], desative "Instalar apps desconhecidos".
  \item \textbf{OPPO/vivo}: [Ajustes] > [Segurança] >\textbf{[Instalar apps de fontes externas]}, desative as autorizações uma a uma.
\end{itemize}

\booksubsection{Mecanismo de verificação de segurança das lojas de apps}

A loja de apps não é “apenas decorativa”. Por trás dela há um pipeline completo de segurança:

\begin{itemize}
  \item \textbf{Verificação de vírus}: Cada app é analisado por vários motores antivírus antes da publicação, identificando assinaturas conhecidas de trojans e vírus.
  \item \textbf{Revisão de permissões}: A loja analisa se as permissões solicitadas pelo app são "razoáveis" — um app de lanterna pedindo acesso aos contatos? Rejeitado ou removido diretamente.
  \item \textbf{Certificação de desenvolvedor}: A publicação requer credenciais empresariais, documento de identidade ou verificação de cartão de crédito. Se houver problemas, é possível identificar o responsável — não são "apps anônimos".
  \item \textbf{Monitoramento contínuo}: Apps já publicados são continuamente reanalisados. Comportamentos maliciosos recém-descobertos resultam em remoção forçada.
\end{itemize}

\booksection{III. Fontes de Terceiros e “Parasitas Eletrônicos”: Risco altíssimo!}

Baixar apps de páginas web ou links suspeitos pode trazer vírus, roubar privacidade e drenar pagamentos. Apresentamos aqui um novo conceito: os\textbf{“Parasitas Eletrônicos”}.

\textbf{O que são “Parasitas Eletrônicos”?}\par

Referem-se a\textbf{softwares “crackeados” com cavalos de Troia embutidos}— você pensa que “economizou dinheiro”, mas na verdade “foi infectado”. Formas comuns:

\begin{itemize}
  \item Softwares pagos em versão "crackeada", "gratuita" ou "sem anúncios";
  \item Ferramentas obscuras como "ganhar curtidas", "roubar envelopes vermelhos", "cheats para jogos";
  \item Os chamados apps "versão interna" recomendados em "anúncios" de grupos de bate-papo.
\end{itemize}

Esses programas parecem funcionar, mas nos bastidores podem estar\textbf{roubando contatos, lendo SMS, escutando o microfone, minerando criptomoedas em segundo plano}. O Modo Puro é a\textbf{primeira linha de defesa}— simplesmente impede que você os instale.

\booksubsection{Sugestões para a segunda linha de defesa:}

\begin{itemize}
  \item \textbf{Verificação completa de vírus periodicamente}: Use o centro de segurança integrado do telefone (como Huawei Optimizer, Xiaomi Security Center) uma vez por mês.
  \item \textbf{Conceda permissões sensíveis com cautela}: Contatos, SMS, localização, câmera, microfone — estas cinco "permissões de alto risco" devem ser rigorosamente aprovadas.
  \item \textbf{Fique atento a consumo anormal de bateria e dados}: Telefone esquentando de repente, bateria descarregando rápido, consumo de dados durante a madrugada — muito provavelmente há um "parasita" escondido em segundo plano.
\end{itemize}

Se for absolutamente necessário instalar um app de terceiros:

\begin{enumerate}
  \item Certifique-se de que a fonte é absolutamente confiável.
  \item Antes de instalar, examine cuidadosamente as permissões solicitadas.
  \item Após a instalação, faça uma verificação com um app de segurança.
\end{enumerate}

Recomendação definitiva: nas configurações do sistema, desative a permissão “Instalar apps desconhecidos”.

\booksection{IV. Instalar, Atualizar, Desinstalar}

\begin{itemize}
  \item \textbf{Instalar}: Na loja, toque em "Obter" ou "Instalar". Após a instalação, na primeira abertura, o app solicitará permissões.
  \item \textbf{Atualizar}:\textbf{Atualize sempre!}Principalmente para corrigir vulnerabilidades de segurança. Recomenda-se ativar a atualização automática.
  \item \textbf{Desinstalar}: Pressione e segure o ícone e arraste para a área de desinstalação. Para uma limpeza mais completa, vá em [Ajustes] > [Gerenciamento de apps] e desinstale por lá.
\end{itemize}

\booksection{V. Gerenciamento de Permissões: Você é o “Chefe” dos Apps}

Apps pedem permissões (como contatos, localização). Às vezes é razoável (mapa precisa de localização), às vezes não (lanterna pedindo contatos).

\begin{itemize}
  \item \textbf{Opções de permissão}: "Permitir sempre", "Permitir apenas durante o uso", "Negar", "Perguntar sempre".
  \item \textbf{Princípios de gestão}:
  \begin{enumerate}
    \item Minimização: não conceda o que não for necessário.
    \item Auditoria posterior: verifique periodicamente nas configurações de permissões quem está "extrapolando".
    \item Contextualização: como apps de transporte, dê permissão de localização apenas durante o uso.
  \end{enumerate}
\end{itemize}

\begin{tcbwarn}{Permissões que merecem atenção especial}
\begin{itemize}
  \item Contatos/SMS: podem ser usados para enviar spam.
  \item Localização em tempo real: pode ser usada para rastrear seus passos.
  \item Câmera/Microfone: podem ser usados para gravação clandestina.
  \item Acesso a arquivos: pode ler seus arquivos privados.
\end{itemize}

\end{tcbwarn}

\booksection{VI. Componentes de um Pacote de Instalação de Software}

Você já se perguntou o que há dentro do pacote de instalação baixado da loja de apps? Entender a composição do pacote ajuda a compreender melhor como os apps funcionam e a saber por onde investigar quando surgirem problemas.

\booksubsection{Estrutura do arquivo APK Android}

O APK (Android Package Kit) é essencialmente um arquivo ZIP. Ao descompactá-lo, você verá os seguintes diretórios e arquivos principais:

\begin{itemize}
  \item \textbf{AndroidManifest.xml}: O "documento de identidade" e "manual" do aplicativo. Registra o nome do pacote, número da versão, permissões necessárias, versão do sistema suportada, Activity de entrada e outras informações essenciais. O sistema lê este arquivo primeiro ao instalar o app.
  \item \textbf{classes.dex}: O código-fonte do aplicativo. O código Java/Kotlin é compilado em bytecode Dalvik e armazenado neste arquivo. Um app pode ter múltiplos arquivos dex (como classes2.dex, classes3.dex).
  \item \textbf{resources.arsc}: Tabela de índice de recursos compilados. Registra todos os IDs de recursos e seus mapeamentos.
  \item \textbf{Diretório res/}: Pasta de recursos. Contém imagens, arquivos de layout, strings, definições de cores, etc. Imagens de diferentes resoluções são colocadas em pastas drawable diferentes (como drawable-hdpi, drawable-xhdpi).
  \item \textbf{Diretório assets/}: Pasta de recursos brutos. Armazena recursos não compilados, como arquivos web, fontes, áudio, etc. O app pode lê-los em tempo de execução via AssetManager.
  \item \textbf{Diretório lib/}: Pasta de bibliotecas nativas. Armazena arquivos .so compilados (produto da compilação de código C/C++), classificados por arquitetura de CPU (como arm64-v8a, armeabi-v7a). Jogos e apps que exigem computação de alto desempenho geralmente têm este diretório.
  \item \textbf{Diretório META-INF/}: Pasta de informações de assinatura. Contém três arquivos: MANIFEST.MF, CERT.SF e CERT.RSA, usados para verificar a integridade e a assinatura do APK.
\end{itemize}

\booksubsection{Estrutura do arquivo IPA do iOS}

O IPA (iOS App Store Package) também é um arquivo compactado, com estrutura ligeiramente diferente:

\begin{itemize}
  \item \textbf{Diretório Payload/}: Onde fica o conteúdo principal. Contém uma pasta .app com todo o conteúdo do aplicativo.
  \item \textbf{Info.plist}: Arquivo de configuração do aplicativo, similar ao AndroidManifest.xml do Android. Registra o ID do app, versão, permissões necessárias, dispositivos suportados, etc.
  \item \textbf{Arquivo executável}: Programa principal com o mesmo nome da pasta .app, é um arquivo binário compilado no formato Mach-O.
  \item \textbf{Diretório Frameworks/}: Bibliotecas dinâmicas e frameworks dos quais o app depende.
  \item \textbf{Arquivos de recursos}: Imagens, áudio, vídeo e outros recursos, normalmente armazenados dentro da pasta .app.
  \item \textbf{Diretório \_CodeSignature/}: Informações de assinatura de código, usadas para verificar a integridade e a origem do aplicativo.
\end{itemize}

\booksubsection{Mecanismo de verificação de assinatura}

Por que os apps baixados de lojas oficiais são mais seguros? A verificação de assinatura é a chave.

\begin{itemize}
  \item \textbf{Assinatura Android}: O desenvolvedor assina o APK com uma chave privada. O sistema verifica com a chave pública durante a instalação. Se o APK for adulterado (por exemplo, com injeção de vírus), a assinatura se torna inválida e o sistema recusa a instalação. O Google Play também exige o App Signing, com o Google armazenando a chave de assinatura, protegendo ainda mais a segurança do app.
  \item \textbf{Assinatura iOS}: A revisão da Apple é mais rigorosa. Cada IPA deve ser assinado com um certificado emitido pela Apple, caso contrário não pode ser instalado. É por isso que o iOS só permite instalar apps da App Store ou via certificado empresarial. Apps assinados com certificado empresarial têm validade limitada e a Apple pode revogá-los a qualquer momento.
  \item \textbf{Função da assinatura}:
  \begin{enumerate}
    \item Verificar a origem: confirmar que o app realmente vem do desenvolvedor declarado.
    \item Verificar a integridade: garantir que o app não foi adulterado durante a transmissão.
    \item Prevenir falsificação: apps maliciosos não podem se passar por outros apps.
  \end{enumerate}
\end{itemize}

\booksubsection{Tamanho do pacote de instalação e otimização}

Você pode notar que o tamanho do app exibido na loja é diferente do espaço ocupado após a instalação. Isso ocorre porque:

\begin{itemize}
  \item \textbf{Tamanho de download}: Tamanho do APK/IPA compactado.
  \item \textbf{Tamanho de instalação}: Tamanho após descompactação, geralmente muito maior que o tamanho de download.
  \item \textbf{Ocupação em tempo de execução}: Durante a execução, o app também gera cache e arquivos de dados, ocupando ainda mais espaço.
\end{itemize}

Técnicas de otimização dos apps modernos:

\begin{itemize}
  \item \textbf{App Bundle (Android)}: Tecnologia de distribuição dinâmica do Google Play. Ao baixar, o usuário obtém apenas os recursos adequados ao seu dispositivo (como baixar apenas as imagens da resolução de tela correspondente), reduzindo significativamente o tamanho do download.
  \item \textbf{App Thinning (iOS)}: Tecnologia similar. A Apple recorta automaticamente os recursos desnecessários de acordo com o dispositivo.
  \item \textbf{Download sob demanda}: Apps grandes como jogos podem baixar primeiro o pacote principal e depois os pacotes de recursos adicionais.
\end{itemize}

\begin{tcbinfo}{Curiosidade: Por que alguns apps têm pacotes de instalação tão grandes?}
Apps de jogos costumam ser os maiores, pois contêm muitas texturas em alta definição, modelos 3D e arquivos de áudio. Apps de redes sociais também são relativamente grandes, pois precisam suportar vários módulos funcionais. Apps utilitários costumam ser os menores. Se você encontrar um app utilitário particularmente grande, desconfie: pode estar empacotado com funcionalidades desnecessárias.

\end{tcbinfo}

\begin{tcbpractice}{Prática: Audite as permissões de 3 apps}
Escolha 3 apps comuns (como redes sociais, compras, utilitários), vá até a página de gerenciamento de permissões em [Ajustes], revise item por item e desative ou altere para “Permitir apenas durante o uso” as permissões sensíveis desnecessárias.

\end{tcbpractice}
